\section{Case-Analysis Details}
\label{app:case-analysis-details}

\subsection{Algorithm~1}
\label{app:case-analysis-details-algo-1}

\subsubsection{Algorithm~1, Case 1(a)}
\label{app:case-one-low}

We show below that it is a \(1/2\)-WSNE:

\begin{itemize}
  \item Row player: since \((\mathbf{x}^*, \mathbf{y}^*)\) is a NE of \((\mathcal{R}, -\mathcal{R})\), for every pure strategy \(i\), \(\mathcal{R}(e_i, \mathbf{y}^*) \le \mathcal{R}(\mathbf{x}^*, \mathbf{y}^*) \le 1/2\). Hence the row player's best-response payoff is \(\le 1/2\), so every row strategy is a \(1/2\)-best response to \(\mathbf{y}^*\);
  \item Column player: since \((\hat{\mathbf{x}}, \hat{\mathbf{y}})\) is a NE of \((-\mathcal{C}, \mathcal{C})\) and the case condition gives \(\mathcal{C}(\hat{\mathbf{x}}, \hat{\mathbf{y}}) \le \mathcal{R}(\mathbf{x}^*, \mathbf{y}^*) \le 1/2\), for every pure strategy \(j\), \(\mathcal{C}(\hat{\mathbf{x}}, e_j) \le \mathcal{C}(\hat{\mathbf{x}}, \hat{\mathbf{y}}) \le \mathcal{R}(\mathbf{x}^*, \mathbf{y}^*) \le 1/2\), so every column strategy is a \(1/2\)-best response to \(\hat{\mathbf{x}}\).
\end{itemize}

Therefore, the strategy profile \((\hat{\mathbf{x}}, \mathbf{y}^*)\) is a \(1/2\)-WSNE.

\subsubsection{Algorithm~1, Case 1(b)}
\label{app:case-one-low-high}

We now show that \((\mathbf{x}', \mathbf{y}^*)\) is a \(1/2\)-WSNE:

\begin{itemize}
  \item Row player: since \(\supp(\mathbf{x}') \subseteq \supp(\mathbf{x}^*)\) and every strategy in \(\supp(\mathbf{x}^*)\) is a best response to \(\mathbf{y}^*\), when the row player uses \(\mathbf{x}'\), every pure strategy in the support is an exact best response, with payoff equal to \(\mathcal{R}(\mathbf{x}^*, \mathbf{y}^*) > 1/2\). That is, the row player's strategy is a \(0\)-best response.
  \item Column player: by construction, \(\mathcal{C}(\mathbf{x}', e_j) \le 1/2\) for all \(j\), so the column player's best-response payoff is \(\le 1/2\), and hence every column strategy is a \(1/2\)-best response to \(\mathbf{x}'\).
\end{itemize}

\subsubsection{Algorithm~1, Case 1(c)}
\label{app:case-one-both-high}

The proof is by contradiction:

\begin{itemize}
  \item If the row player's payoff \(\le 1/2\): since \(\mathbf{x}^*\) is a maximin strategy, \(\mathcal{R}(\mathbf{x}^*, \mathbf{y}) \ge \mathcal{R}(\mathbf{x}^*, \mathbf{y}^*) > 1/2\), so he can switch to \(\mathbf{x}^*\) to improve his payoff, a contradiction;

  \item If the column player's payoff \(\le 1/2\):

        Consider the zero-sum game \((-\mathcal{C}_{\mathbf{x}^*}, \mathcal{C}_{\mathbf{x}^*})\), and let \((\tilde{\mathbf{x}}, \tilde{\mathbf{y}})\) be a NE of it. Suppose \(\mathcal{C}(\tilde{\mathbf{x}}, \tilde{\mathbf{y}}) \le 1/2\); then
        \[
          \mathcal{C}(\tilde{\mathbf{x}}, e_j) \le \mathcal{C}(\tilde{\mathbf{x}}, \tilde{\mathbf{y}}) \le 1/2 \;\; \forall j \in [n]
        \]
        Moreover, \(\supp(\tilde{\mathbf{x}}) \subseteq \supp(\mathbf{x}^*)\). Thus \(\tilde{\mathbf{x}}\) satisfies the feasibility system in Case 1(b), a contradiction. Hence \(\mathcal{C}(\tilde{\mathbf{x}}, \tilde{\mathbf{y}}) > 1/2\).

        Since \(\tilde{\mathbf{y}}\) is the column player's maximin strategy in the zero-sum game, for every \(\mathbf{x}\) satisfying \(\supp(\mathbf{x}) \subseteq \supp(\mathbf{x}^*)\),
        \[
          \mathcal{C}(\mathbf{x}, \tilde{\mathbf{y}}) \ge \mathcal{C}(\tilde{\mathbf{x}}, \tilde{\mathbf{y}}) > \frac{1}{2}
        \]
        Then the column player can improve his payoff by switching to \(\tilde{\mathbf{y}}\), contradicting that \((\mathbf{x}, \mathbf{y})\) is a NE.
\end{itemize}

\subsection{Algorithm~2}
\label{app:case-analysis-details-algo-2}

\subsubsection{Algorithm~2, Case 2(a)}
\label{app:case-two-low}

\begin{itemize}
  \item Row player: \(1/2 \ge \mathcal{R}(\mathbf{x}^*, \mathbf{y}^*) \ge \mathcal{R}(e_i, \mathbf{y}^*) - \varepsilon\), so every strategy is a \((1/2 + \varepsilon)\)-best response;
  \item Column player: \(1/2 \ge \mathcal{R}(\mathbf{x}^*, \mathbf{y}^*) \ge \mathcal{C}(\hat{\mathbf{x}}, \hat{\mathbf{y}}) \ge \mathcal{C}(\hat{\mathbf{x}}, e_j) - \varepsilon\) (by Observation 1), so every strategy is a \((1/2 + \varepsilon)\)-best response.
\end{itemize}

Hence \((\hat{\mathbf{x}}, \mathbf{y}^*)\) is a \((1/2 + \varepsilon)\)-WSNE, and in particular a \((1/2 + 3\varepsilon + \delta)\)-WSNE.

\subsubsection{Algorithm~2, Case 2(b)}
\label{app:case-two-low-high}

We use the two lemmas to show that \((\mathbf{x}', \mathbf{y}_s^*)\) is a \((1/2 + 3\varepsilon)\)-WSNE:

\begin{itemize}
  \item Row player: by Lemma~\ref{lem:3_eps-wsne-sampling}, \((\mathbf{x}_s^*, \mathbf{y}_s^*)\) is a \(3\varepsilon\)-WSNE, and \(\mathbf{x}'\) is supported on \(\supp(\mathbf{x}_s^*)\), which means that every pure strategy in the support of \(\mathbf{x}'\) is a \(3\varepsilon\)-best response;
  \item Column player: by construction, the column player's maximum payoff is at most \(1/2\).
\end{itemize}

Hence \((\mathbf{x}', \mathbf{y}_s^*)\) is a \((1/2 + 3\varepsilon)\)-WSNE.

\subsubsection{Algorithm~2, Case 2(c)}
\label{app:case-two-both-high}

The proof is again by contradiction:

\begin{itemize}
  \item Suppose the row player's payoff is at most \(1/2-3\varepsilon\). Applied to the column payoff \(-\mathcal{R}\) in the zero-sum game, Lemma~\ref{lem:small-supp-approx} gives
        \(|\mathcal{R}(\mathbf{x}^*,\mathbf{y})-\mathcal{R}(\mathbf{x}_s^*,\mathbf{y})|\le\varepsilon\) for every mixed strategy \(\mathbf{y}\). Observation 1 also gives
        \(\mathcal{R}(\mathbf{x}^*,\mathbf{y})\ge\mathcal{R}(\mathbf{x}^*,\mathbf{y}^*)-\varepsilon\).
        Consequently,
        \[
          \mathcal{R}(\mathbf{x}_s^*, \mathbf{y}) \ge \mathcal{R}(\mathbf{x}^*, \mathbf{y}) - \varepsilon \ge \mathcal{R}(\mathbf{x}^*, \mathbf{y}^*) - 2\varepsilon > \frac{1}{2} - 2\varepsilon
        \]
        and the row player can profitably deviate to \(\mathbf{x}_s^*\), a contradiction.

  \item Suppose the column player's payoff is at most \(1/2-3\varepsilon\). The maximin argument from Case 1(c) applies to the restricted game \((-\mathcal{C}_{\mathbf{x}_s^*},\mathcal{C}_{\mathbf{x}_s^*})\). Indeed, if its equilibrium value were at most \(1/2\), its row strategy would satisfy the feasibility condition of Case 2(b), contrary to the premise of this case. Thus its column maximin strategy \(\tilde{\mathbf{y}}\) satisfies, for every row strategy \(\mathbf{x}\) satisfying \(\supp(\mathbf{x}) \subseteq \supp(\mathbf{x}_s^*)\),
        \[
          \mathcal{C}(\mathbf{x},\tilde{\mathbf{y}})>\frac12>\frac12-3\varepsilon.
        \]
        This is a profitable deviation, again a contradiction.
\end{itemize}

\subsection{The Sampled-Profile Guarantee}
\label{app:sampled-wsne-proof}

\textbf{Proof:} For any \(i \in \supp(\mathbf{x}_s) \subseteq \supp(\mathbf{x})\), Lemma~\ref{lem:small-supp-approx} and the \(\varepsilon\)-WSNE property give
\begin{align*}
  \mathcal{R}(e_i, \mathbf{y}_s)
   & \ge \mathcal{R}(e_i, \mathbf{y})-\varepsilon                  \\
   & \ge \max_{k\in[n]}\mathcal{R}(e_k,\mathbf{y})-2\varepsilon    \\
   & \ge \max_{k\in[n]}\mathcal{R}(e_k,\mathbf{y}_s)-3\varepsilon.
\end{align*}
The column-player inequality is symmetric, so \((\mathbf{x}_s,\mathbf{y}_s)\) is a \(3\varepsilon\)-WSNE.
